% Table: C-MAPSS benchmark evaluation (shared-model regime, quality-controlled)
\begin{table}[t]
\centering
\begin{threeparttable}
\caption{\textbf{Cross-view energy distances on the C-MAPSS FD001 simulated
turbofan benchmark under the shared-model protocol.} A single model per
architecture is trained on data that mixes all three observation
configurations, mirroring the synthetic protocol. Lower energy distance
indicates closer agreement between dropout-sampled predictive distributions
for the same engine viewed under different sensor configurations; distances
are on the raw RUL scale (cycles). Bootstrap 90\% confidence intervals in
brackets (2,000 replicates, resampling engines). Negative relative reduction
indicates PROP-DIRECT is \emph{worse}.}
\label{tab:cmapss_validation}
\begin{tabular}{lccc}
\toprule
\textbf{View Pair} & \textbf{PROP-DIRECT} & \textbf{B-STDPROB} & \textbf{Relative Reduction} \\
\midrule
Full vs.\ View 1    & 1.85 [1.61, 2.11] & 2.27 [2.04, 2.51] & 18.5\% [7.8\%, 29.3\%] \\
Full vs.\ View 2    & 1.80 [1.56, 2.05] & 2.54 [2.29, 2.79] & 28.9\% [18.4\%, 37.9\%] \\
View 1 vs.\ View 2  & 1.85 [1.57, 2.17] & 1.36 [1.19, 1.52] & $-36.6$\% [$-62.9$\%, $-13.1$\%] \\
\bottomrule
\end{tabular}
\begin{tablenotes}
\footnotesize
\item Single-view test MAE (cycles), same runs: PROP-DIRECT 16.44 (full) /
15.95 (View 1) / 16.25 (View 2); B-STDPROB 16.27 (full) / 15.42 (View 1) /
15.73 (View 2). The near-identical point accuracy makes this geometry
comparison quality-controlled: differences in cross-view distance cannot be
attributed to differences in posterior informativeness.
\end{tablenotes}
\end{threeparttable}
\end{table}
